\section{Discussion}

\subsection{Paired correspondence makes scanner suppression testable}
The principal result is not a lower scanner-probe score by itself. It is the simultaneous reduction in tissue-branch scanner recoverability, preservation of same-region retrieval, and greater scanner accessibility in the acquisition branch relative to an equal-capacity neural control. These three observations address different failure modes: a scanner-confused representation can collapse or lose tissue identity, while a reconstructive representation can retain scanner information indiscriminately. PA-NF supplies an explicit paired-region constraint against both outcomes.

The pairing itself matters. In broken-correspondence controls, invalid region assignments can reduce scanner-probe accuracy while damaging same-region geometry. The favorable result with genuine pairs therefore cannot be explained by scanner confusion alone. Cross-backbone tests show that the objective can operate on different frozen feature families; the independent canine study evaluates a different, harder frontier against strong linear scanner corrections. Together, these studies define the scope of the original controlled result without replacing its registered comparator.

\subsection{External retrieval without retraining}
Two independently acquired ScanComp slides provide a distinct test of the trained representation. The 25 original PA-NF checkpoints were applied without retraining or per-slide checkpoint selection. Strict cross-scanner top-1 retrieval improved by 4.53 and 6.97 percentage points over raw DINOv2. Both slides also had positive average margin contrasts against the matched no-adversary and no-scanner-objective control families. These observations support external technical transfer of region discrimination, rather than only an in-cohort scanner-probe difference.

Registration is part of the evaluation boundary. Both reported panels qualified for the three-scanner analysis under unchanged geometric rules; neither satisfied full four-scanner admission. Other source-selected slides that failed registration were not run through the neural models and remain part of the eligibility record. Each ScanComp slide, not each checkpoint or query, is the relevant physical specimen. The two positive panels therefore support a descriptive replication of retrieval, not a population-level effect estimate or a diagnostic-label conclusion.

\subsection{Source influence is a second, separately tested problem}
The downstream stages ask how a system should aggregate information once representations or slide-level predictions exist. TransnnMIL and established MIL baselines test slide aggregation on PANDA under matched conditions. The institutional experiments test a different decision: whether assigning influence by source volume remains appropriate after validation behavior changes. On simulated PANDA sites, a clean-calibrated detector switches aggregation under dominant-site perturbation; in CAMELYON17, changing source weights improves performance at a held-out center. Both results are informative about institutional influence, but neither measures the causal effect of transferring PA-NF features through a jointly trained multi-site model.

\subsection{Interpretation}
PA-NF's methodological advance is a correspondence-aware training and evaluation strategy in which acquisition accessibility is constrained subject to measured preservation of tissue-region identity. Its strongest evidence is the registered SCORPION neural-control comparison, strengthened by two frozen-checkpoint external retrieval observations. Full biological--acquisition identifiability is not established by these metrics. A larger independently qualified human cohort with authenticated biological labels would determine whether the observed representation advantage extends to clinically meaningful information.
